\documentclass[10pt,a4paper]{amsart}
\usepackage[margin=19mm]{geometry}
\usepackage{amsmath,amssymb,mathtools}
\usepackage[scr=rsfs,cal=euler]{mathalfa}
\usepackage{xfrac}
\usepackage[unicode,hidelinks,bookmarks=false]{hyperref}
\newcommand{\ie}{i.e.\ }
\newcommand{\cf}{cf.\ }
\newcommand{\resp}{respectively\ }
\newcommand{\Subsection}{\subsection}
\providecommand{\nobreakdash}{\nobreak\hbox{-}}
\setlength{\emergencystretch}{2em}
\multlinegap0pt
\usepackage{bm}
\usepackage{bbm}
\usepackage{dsfont}
\usepackage{xspace}
\usepackage{cancel}
\usepackage{mathtools}
\usepackage[shortlabels]{enumitem}
\usepackage{amsthm}
\makeatletter
\@ifundefined{Theorem}{\newtheorem{Theorem}{Theorem}}{}
\@ifundefined{Corollary}{\newtheorem{Corollary}[Theorem]{Corollary}}{}
\@ifundefined{Lemma}{\newtheorem{Lemma}[Theorem]{Lemma}}{}
\@ifundefined{Proposition}{\newtheorem{Proposition}[Theorem]{Proposition}}{}
\makeatother
\def\Iset{{\mathbb I}}
\def\deg{\operatorname{deg}}
\def\reg{\mathrm{reg}} 
\def\wbf{{\mathbf w}}
\title[New bounds on Castelnuovo--Mumford regularity of monomial cu]{New bounds on Castelnuovo--Mumford regularity of monomial curves and application to sumsets}
\author{Le Tuan Hoa, Doan Quang Tien}
\date{Source: arXiv:2604.20189v1 (CC BY 4.0)}
\begin{document}
\maketitle

\noindent\emph{Source and licence.} New bounds on Castelnuovo--Mumford regularity of monomial curves and application to sumsets. Version arXiv:2604.20189v1 (CC BY 4.0), distributed under the Creative Commons Attribution 4.0 International licence, as recorded in the arXiv licence metadata for this exact version and shown on the abstract page https://arxiv.org/abs/2604.20189v1. Full rights notice: licenses/arxiv-2604.20189-CC-BY-4.0.txt.\par\smallskip

\noindent\textit{Editorial scope.} Counted unit: the Proposition whose source label is \texttt{Bd2} in the article (source0.tex lines 617--619, source label \texttt{Bd2}), delivered together with the article's own complete original proof of it (source0.tex lines 621--687, one \texttt{proof} environment of 67 lines). No mathematics is written, repaired or re-derived: statement and proof are the article's own text line for line. Required context retained uncounted: B2 (lines 414--415); A4 (lines 485--494); A5 (lines 533--541). Every definition, display and label of those lines is retained unchanged. Omitted from this package and not redistributed: 1--413, 416--484, 495--532, 542--616, 620--620, 688--1202. Nothing inside a retained excerpt is omitted or altered: every retained body is delivered exactly as the article's lines are written, so no omission receipt of any kind accompanies this package. Citations: the retained excerpts carry no citation command at all, so no bibliography entry of the article is transcribed and no reference is redistributed. Reference adaptation: none was needed and none was made; every reference inside the delivered unit resolves inside it, so no unresolved reference or citation is produced. Printed slips kept literal: no printed word, symbol, hypothesis or number of the counted statement or of its proof was altered. Layout and class: the article's own class is \texttt{article} with packages including amsthm, amsfonts, amssymb, amsmath, mathrsfs, enumerate, hyperref, geometry; the wrapper is the lane's pinned \texttt{amsart} editorial wrapper at 10pt on A4, which restates the theorem-like environments the retained text needs and adds the article's own macro definitions that the retained text needs (\texttt{\textbackslash Iset}, \texttt{\textbackslash deg}, \texttt{\textbackslash reg}, \texttt{\textbackslash wbf}), with extra wrapper packages (bm, bbm, dsfont, xspace, cancel, mathtools, enumitem). The delivered numbering is the unit's own. Assets: no retained span contains a figure environment, a graphics-inclusion command, a PDF-page inclusion, a table or a TikZ picture; no image file is copied into this package, the package declares no asset dependency and no image file is redistributed with it. Licence: version arXiv:2604.20189v1 of this article is distributed under the Creative Commons Attribution 4.0 International licence, as recorded in the arXiv licence metadata for this exact version and shown on the abstract page https://arxiv.org/abs/2604.20189v1. A full rights notice is delivered at licenses/arxiv-2604.20189-CC-BY-4.0.txt. Characters: every retained excerpt is pure ASCII, so the delivered engine, XeLaTeX with its default Latin Modern text font, reports no missing character.
\begin{Lemma} \label{B2}  $a_2(K[S]) \le \left\lfloor \frac{F_1 + F_2}{d} \right\rfloor$.
\end{Lemma}

\begin{Theorem} \label{A4} Assume that $\Iset \neq \emptyset$. Then
\begin{enumerate}[\((i)\)]
    \item $a_1(K[S])  \geq \max_{i\in \Iset}
\{\delta_1(\wbf_i),\ \delta_2(\wbf_i)\} - 1$;
\item $a_1(K[S])  \leq \max_{i\in \Iset}
 \{\delta_1(\wbf_i) +  \delta_2(\wbf_i) - \deg(\wbf_i) \} - 2 $;
 \item $a_1(K[S])  \leq  \left\lfloor \left( 1- \frac{a_1}{d} \right)  \max_{i\in \Iset} \{\delta_1(\wbf_i) \} +  \frac{a_{k-1}}{d}  \max_{i\in \Iset}
\{  \delta_2(\wbf_i) \} \right\rfloor - 2$.
\end{enumerate}
\end{Theorem}

\begin{Corollary} \label{A5} 
\begin{enumerate}[\((i)\)]
    \item Assume that $a_1=1$ and $a_{k-1} = d-1$. Let $\varepsilon = \max\{i|\ 1,\ldots,i, d-1,\ldots,d-i \in \{a_1,\ldots,a_{k-1} \}\}$ and $\lambda_{\max}$ be the largest gap of $A_1$. Then
$$\reg(K[S]) \leq \left\lfloor \frac{\lambda_{\max} -1}{\varepsilon} \right\rfloor + 2.$$
\item Assume that $a_1=1,\ldots,a_p = p, a_{p+1} \ge p+2$ for some $p\ge 1$. Denote the first non-zero gap of $A_1$ by $\lambda = a_{p+1}-p-1$. Then
$$\reg(K[S]) \geq \left\lceil \frac{\lambda}{p} \right\rceil + 1 = \left\lfloor \frac{\lambda -1}{p} \right\rfloor + 2.$$
In particular,  the bound in \((i)\) is attained if $\lambda = \lambda_{\max}$ and $ \varepsilon = p$.
\end{enumerate}
\end{Corollary}

\begin{Proposition} \label{Bd2} Let $A_1 = \{a_0,\ldots,a_k\} = \bigsqcup_{j=0}^r [b_{2j},b_{2j+1}]$, where $0=:b_0 \leq b_1 < b_2 \leq b_3 < \cdots < b_{2r}\leq b_{2r+1} := d$ and $b_{2j-1} + 2 \leq b_{2j}$ for $j=1,\ldots,r$, $r\ge 2$.  If $b_1 = 0, b_{2r} < d$ and $b_5 - b_4 + 1\geq b_2$, then
$$\reg(K[S]) \leq 2+ \max\left\{ \left\lfloor \frac{ \lambda_{\mathrm{sl}}}{d-b_{2r}} \right\rfloor + \left\lfloor \frac{\lambda_{\max}}{b_2} \right\rfloor;\ \left\lfloor \frac{ \lambda_{\max}}{d-b_{2r}} \right\rfloor + \left\lfloor \frac{\lambda_{\mathrm{sl}}}{b_2} \right\rfloor \right\}.$$
\end{Proposition} 

\begin{proof} Let
$$M:=  \max\left\{ \left\lfloor \frac{ \lambda_{\mathrm{sl}}}{d-b_{2r}} \right\rfloor + \left\lfloor \frac{\lambda_{\max}}{b_2} \right\rfloor;\ \left\lfloor \frac{ \lambda_{\max}}{d-b_{2r}} \right\rfloor + \left\lfloor \frac{\lambda_{\mathrm{sl}}}{b_2} \right\rfloor \right\}\geq 0.$$
From the assumption we see that $1\in A_2$. Hence $F_2 = -1$ and $\omega_2(d- x) = d- x$ for all $x\leq d$. Since $b_5 - b_4 + 1\geq b_2$, $b_1= 0$,  we have $1 \leq F_1 \leq b_4 -1$ and
$$\omega_1(x) = \begin{cases} x & \text{if}\quad  b_4 \leq x < d \ \text{or} \ x \in \langle [b_2,b_3] \rangle \ \text{and}\ x< b_4,\\
x+ d &\text{otherwise}.
\end{cases}$$
 By Lemma \ref{B2}, $a_2(K[S]) + 2 \leq \lfloor \frac{F_1+ F_2}{d}\rfloor + 2 = 2$. So, it is left to  show that $a_1(K[S]) \leq M + 1 $, provided  that $\Iset\neq \emptyset$. By Theorem \ref{A4}\((ii)\) there is  $x\in \Iset$ such that
 $$a_1(K[S]) \leq \delta_1(\wbf_x ) + \delta_2(\wbf_x) - \deg(\wbf_x) - 2 =: D(x).$$
 It is enough to show that $D(x) \leq M + 1$. Let $\epsilon' := \min\{ b_2 , d-b_{2r}\}$. We distinguish three cases.
 \begin{enumerate}[\(\textbf{\text{Case}}\ \)\bf a.]
     \item $ x>b_5$. In this case $\wbf_x = (x, d-x)$, whence $\deg(\wbf_x) = 1$.  There is $i$ such that $b_5\leq a_i < x < a_{i+1}$. In the proof of Corollary \ref{A5}\((i)\) we know that 
 \begin{equation}\label{EBd21}
\delta_2(\wbf_x) \leq 1 + \left\lceil \frac{a_{i+1} - x}{d- b_{2r}} \right\rceil \leq 1+ \left\lceil \frac{a_{i+1} - x}{\epsilon'} \right\rceil.
\end{equation}
  Note that $x = b_4 + (x-b_4)$. If $x-b_4 \in A_1$, then $\delta_1(\wbf_x) = \delta_1(x) = 2$. We now assume that $x-b_4 \not\in A_1$
  \begin{enumerate}[\((\text{a}1)\)]
      \item If $ x-b_4 > a_i$, then using $x= b_4 + x-b_4 = (b_4 + b'_2) + a_i +  \left\lfloor \frac{x-b_4 - a_i}{b_2} \right\rfloor b_2$, where $0\leq b'_2 < b_2$, and noticing that $b_4 + b'_2 \leq b_5$, whence $b_4+b'_2 \in A_1$, we can conclude that 
  \begin{equation}\label{EBd22}
  \delta_1(\wbf_x) = \delta_1(x) \leq 2 + \left\lfloor \frac{x-b_4 - a_i}{b_2} \right\rfloor  \leq 2+ \left\lfloor \frac{x-b_4 - a_i}{\epsilon'} \right\rfloor.
  \end{equation}
 From (\ref{EBd21}) and (\ref{EBd22}) we get
  \begin{equation}\label{EBd23} 
  D(x)  \leq \left\lceil \frac{a_{i+1} - x}{\epsilon'} \right\rceil +  \left\lfloor \frac{x-b_4 - a_i}{\epsilon'} \right\rfloor \leq \left\lceil \frac{\lambda_i -3 }{\epsilon'} \right\rceil  \leq 1 + \left\lfloor \frac{\lambda_i  -4}{\epsilon'} \right\rfloor \leq M+1.
   \end{equation}
      \item $a_j < x-b_4 < a_{j+1}$ for some $j<i$, that is $x$ and $x-b_4$ belong to two different gaps of $A_1$. Using $x= b_4 + x-b_4 = (b_4 + b^{''}_2) + a_j +  \left\lfloor \frac{x-b_4 - a_j}{b_2} \right\rfloor b_2$, where $0\leq b^{''}_2 < b_2$, we get
  \begin{equation}\label{EBd24}
  \delta_1(\wbf_x) = \delta_1(x) \leq 2 + \left\lfloor \frac{x-b_4 - a_j}{b_2} \right\rfloor \leq 2 + \left\lfloor \frac{a_{j+1} - 1 - a_j}{b_2} \right\rfloor \leq 2+ \left\lfloor \frac{\lambda_j}{b_2} \right\rfloor.
   \end{equation}
   From (\ref{EBd21}) we also have $\delta_2(\wbf_x) \leq 1 + \left\lceil \frac{ \lambda_i}{d-b_{2r}} \right\rceil$. Hence,
    \begin{equation}\label{EBd25}
  D(x)  \leq \left\lceil \frac{ \lambda_i}{d-b_{2r}} \right\rceil + \left\lfloor \frac{\lambda_j}{b_2} \right\rfloor \leq 1+  \left\lfloor \frac{ \lambda_i}{d-b_{2r}} \right\rfloor + \left\lfloor \frac{\lambda_j}{b_2} \right\rfloor.
     \end{equation}
   Note that if $\lambda_k\ge \lambda_l$ are two different gaps, then $\lambda_k \leq \lambda_{\max}$ and $\lambda_l \leq \lambda_{\mathrm{sl}}$. This implies
   $\left\lfloor \frac{ \lambda_i}{d-b_{2r}} \right\rfloor + \left\lfloor \frac{\lambda_j}{b_2} \right\rfloor \leq  M.$ Using this remark, from (\ref{EBd23}) and (\ref{EBd25}) we get in Case $a$ that $ D(x) \leq  M +1$.
  \end{enumerate}
     \item $ b_3 < x< b_4$.
     \begin{enumerate}[\((\text{b}1)\)]
         \item $x\not\in \langle [b_2,b_3]\rangle$. Then $\omega_1(x) = x+d$, whence $\wbf_x = (d+x, d-x)$ and $\deg(\wbf_x) = 2$. Assume that $b_3 = a_i$. Then $a_{i+1} = b_4$, and $x$ belongs to the gap $\lambda_i$. Write $d+x = b_4 + (d+x-b_4)$. Note that $b_3 <  d+x - b_4 < d$. If $d+x-b_4 \in A_1$, then $ \delta_1(\wbf_x) = \delta_1(d+ x)  =2$. Assume now that $d+x-b_4\not\in A_1$.
  
   If $d+x-b_4$ belongs to $j$-th gap with $i\neq j$, then similar to the subcase \(\text{(a2)}\), we can conclude  that 
  $$\delta_1(\wbf_x) = \delta_1(d+ x)  \leq 2+ \left\lfloor \frac{\lambda_j}{b_2} \right\rfloor.$$
  On the other hand,  using the relation $d-x = (d-b_4) + (b_4-x)$,  as shown in the proof of Corollary \ref{A5}\((i)\), 
   \begin{equation}\label{EBd26}
   \delta_2(\wbf_x)  \leq 1 +  \left\lceil \frac{b_4-x}{d-b_{2r}} \right\rceil \leq  1+ \left\lceil \frac{\lambda_i}{d-b_{2r}} \right\rceil \leq 2+ \left\lfloor \frac{\lambda_i -1}{d-b_{2r}} \right\rfloor.
    \end{equation}
   Hence,  the argument at the end of Case \text{(a2)} gives  $D(x) \le M$.
  
  The remaining case is  $b_3< d+x-b_4 < b_4$; that is both $x$ and $d+x-b_4$ belong to the same $i$-th gap. Using the relation $d+x = b_3 + b_4 + (d+x- b_3-b_4)$, similar to the subcase (a1) we get
  $$\delta_1(\wbf_x) = \delta_1(d+ x)  \leq 2+ \left\lfloor \frac{d+x-b_3-b_4}{b_2} \right\rfloor \leq 2+ \left\lfloor \frac{d+x-b_3-b_4}{\epsilon'} \right\rfloor.$$
  On the other hand $d-x = 2(d-b_4) + (b_4 - (d+x-b_4))$. Note that $1, d-b_4 \in A_2$. As shown in the proof of Corollary \ref{A5}\((i)\), 
   $$\delta_2(\wbf_x)  = \delta_2(d-x) \leq 2 + \left\lceil \frac{b_4 - (d+x-b_4)}{d-b_{2r}} \right\rceil \leq 2 +  \left\lceil \frac{2b_4 - (d+x)}{\epsilon'} \right\rceil.$$
   Hence,
   \begin{align*}
       D(x) &\leq \left\lfloor \frac{d+x-b_3-b_4}{\epsilon'} \right\rfloor + 
 \left\lceil \frac{2b_4 - (d+x)}{\epsilon'} \right\rceil \\
&\leq \left\lceil \frac{b_4 - b_3}{\epsilon'} \right\rceil = \left\lfloor \frac{\lambda_i}{\epsilon'} \right\rfloor + 1 \leq M+ 1.
   \end{align*}
         \item $x\in \langle [b_2,b_3]\rangle$. Then $\omega_1(x) =x$ and $\deg(\wbf_x) = 1$. By (\ref{EBd26}),
   $$\delta_2(\wbf_x)  \leq 1+ \left\lceil \frac{b_4 - x}{d-b_{2r}} \right\rceil \leq 1 +  \left\lceil \frac{b_4 - x}{\epsilon'} \right\rceil.$$
   Since $x< b_4$, we can find non-negative integers $n_j$ such that $x =  \sum_{j=b_2}^{b_3} n_jj$ such that $\delta_1(x) =  \sum_{j=b_2}^{b_3} n_j$. Assume that $\delta_1(x) \geq 3$. Let $b$ be a subsum of two summands in the sum $ \sum_{j=b_2}^{b_3} n_jj$ (could be the same). Then $b\geq b_3+1$, because otherwise replacing this subsum by $b$ would give a presentation of $x$ with the sum of coefficients strictly smaller than $\delta_1(x)$, which is impossible. We have $x-b\geq (\delta_1(x) - 2)b_2$. Hence,
   $$\delta_1(\wbf_x) = \delta_1(x) \leq  2+ \left\lfloor \frac{x-b}{b_2} \right\rfloor \leq  2+ \left\lfloor \frac{x-b_3-1}{b_2} \right\rfloor \leq 2+ \left\lfloor \frac{x-b_3-1}{\epsilon'} \right\rfloor.$$
   This implies
   $$D(x) \leq   \left\lceil \frac{b_4 - x}{\epsilon'} \right\rceil + \left\lfloor \frac{x-b_3 -1}{\epsilon'} \right\rfloor \leq  \left\lceil \frac{b_4-b_3 -1}{\epsilon'} \right\rceil = \left\lfloor \frac{\lambda_i}{\epsilon'} \right\rfloor \le M.$$
     \end{enumerate}
     \item $0< x< b_2$, i.e. $x$ belongs to the first gap (note that $a_0 = 0, \ a_1 = b_2$). Then $x\not\in \langle A_1\rangle$, whence $\omega_1(x) =d+x$, $\wbf_x = (d+x, d-x)$ and $\deg(\wbf_x) = 2$. We have $d+ x = b_2 + (d+x-b_2)$. If $d+x-b_2$ belongs to the first gap too, then $d+x-b_2 < b_2$. This implies $d< 2b_2 - 1< b_5$, a contradiction. Hence if $d+x-b_2\not\in A_1$, then $x$ and $d+x-b_2$ belong to two different gaps of $A_1$. The argument in the case \text{(a2)} gives us $D(x)\leq M + 1$, as required. 
 \end{enumerate}
\end{proof}

\end{document}
